\phantomsection

\subsection*{Ingredientes}

\addcontentsline{toc}{subsection}{Ingredientes}

\begin{itemize}

\item Após finalizados, os hidroméis podem receber ajustes de dulçor, acidez ou outras características.
\item Estabilizantes ou outros coadjuvantes de processo não devem deixar sabores perceptíveis no produto final.
\item As variedades de mel variam em intensidade e caráter, sendo algumas mais facilmente reconhecíveis do que outras. Méis de caráter sutil não devem ser penalizados por expressarem suas características próprias. Da mesma forma, méis de caráter mais intenso não são, por si só, superiores aos mais sutis.
\item O uso de madeira relativamente neutra na fermentação ou no envelhecimento é permitido em qualquer estilo, desde que permaneça como caráter de fundo. No entanto, exemplares com caráter de madeira significativo devem ser inscritos em M4D Wood-Aged Mead, enquanto aqueles que apresentem caráter de barril, com sabores provenientes de outras bebidas alcoólicas anteriormente armazenadas nesse barril, devem ser inscritos em M4E Barrel-Aged Mead.
\item As frutas no hidromel podem influenciar os níveis de acidez e taninos, além de contribuir com cor. A fruta não deve parecer artificial, crua ou não fermentada. A cor do hidromel proveniente da fruta geralmente é mais clara que a cor da polpa e pode adquirir tonalidades ligeiramente diferentes após a fermentação. Hidroméis produzidos com frutas de coloração mais clara podem apresentar cor proveniente das variedades de mel utilizadas. Nos hidroméis que formam espuma, ela também pode adquirir parte da coloração da fruta.
\item Ao utilizar especiarias, não é necessário declarar todas elas caso não sejam perceptíveis. Combinações de especiarias amplamente conhecidas podem ser identificadas por seu nome comum, como especiarias para torta de maçã, especiarias para torta de abóbora, curry em pó etc.

\end{itemize}